\documentclass{article}
\usepackage{fullpage,amsmath,amsthm,graphicx,enumitem,amsfonts}
\usepackage{multicol}
\usepackage{booktabs}
\usepackage{algorithm}
\usepackage{algpseudocode}
\usepackage{sgame}

\usepackage{tikz}
\usetikzlibrary{positioning, shapes, arrows, fit}

\theoremstyle{definition}
\newtheorem{thm}{Theorem}
\newtheorem{question}[thm]{Question}
\newenvironment{solution}{\noindent\textit{Solution:}}{}

\title{ASEN 5264 Decision Making under Uncertainty\\
       Exam 3: POMDPs and Games}


\begin{document}
\maketitle


\noindent \textbf{Instructions:} Clearly indicate your final answers and \textbf{briefly justify all answers with text, mathematical expressions, and/or calculations}. If you do not understand how to do a problem, skip it and move on so that you have time to attempt all problems. You may consult any static source, but you may NOT communicate with any person except the instructor or TA, and you may not use LLMs such as ChatGPT.
\\
\\
\noindent\textbf{Important:} Please include the following statement along with your signature (typed or handwritten) on the exam: \textbf{``I affirm that I completed this exam without the use of any artificial intelligence tool such as ChatGPT or Github Copilot.''}

\clearpage

\noindent\textbf{Note:} Questions 1 and 2 have been removed from this exam because they were added to Homework 5 and 6.

\begin{question} (10 points)
    Suppose that Alice and Bob are playing the 5-card version of ``Goofspiel'' or ``Game of Pure Strategy'' discrete Markov game played in class. Alice has computed a Nash equilibrium for the game and she shares her Nash equilibrium policy, $\pi^1_A$, with Bob.
    \begin{enumerate}[label=\alph*)]
        \item If Bob knows the $(I, S, A, T, R, \gamma)$ model of the game and Alice's policy $\pi^1_A$, describe in two or three sentences and equations how Bob can compute a best response policy with a computer.
        \item If Bob then plays the best response policy computed in the manner described above against $\pi^1_A$, is the resulting joint policy guaranteed to be a Nash equilibrium? Briefly justify your answer.
    \end{enumerate}
\end{question}

\end{document}
